
\section{Dashboard Subsystem (DBS)}%
\label{sec:dbs}

Il DBS è l'interfaccia web per il monitoraggio del livello dell'acqua, della modalità operativa, dell'apertura comandata e dello storico, e per l'invio di richieste al CUS.

\subsection{Architettura dell'interfaccia web}

L'interfaccia è sviluppata con HTML, CSS e JavaScript, senza framework frontend.
Comprende pannelli informativi, controlli per l'operatore e un grafico per la visualizzazione dello storico.
Moduli distinti gestiscono la comunicazione HTTP, l'interfaccia e il grafico dello storico, realizzato con Chart.js incluso localmente.

\subsection{Comunicazione HTTP e gestione dello stato}%
\label{subsec:dbs-state}

Il DBS comunica con il CUS tramite HTTP, recuperando periodicamente lo stato del sistema e lo storico delle misure.
Nella configurazione attuale, le richieste sono indirizzate a \url{http://localhost:8080}, con intervalli di aggiornamento nominali di \(1000\,\mathrm{ms}\) per entrambi i dati.

L'interfaccia distingue le misure non disponibili, rappresentate da \code{null}, da un livello valido pari allo \(0\%\).
Nel grafico, le misure non valide sono rappresentate mediante interruzioni della linea.

La dashboard distingue inoltre due condizioni di indisponibilità.
Lo stato \code{UNCONNECTED} viene comunicato dal CUS quando non riceve misure valide dal TMS per un intervallo prolungato.
La condizione \code{NOT AVAILABLE} indica invece che il DBS non riesce a ottenere uno stato utilizzabile dal CUS.

In quest'ultimo caso, i comandi vengono disabilitati e gli ultimi valori disponibili rimangono visualizzati, senza essere considerati informazioni aggiornate.
Un errore nel recupero del solo storico viene invece segnalato separatamente, senza modificare lo stato di disponibilità del CUS.

\subsection{Interazione dell'operatore}

L'interfaccia mette a disposizione un pulsante per richiedere il cambio tra le modalità \code{AUTOMATIC} e \code{MANUAL} e uno slider per regolare l'apertura della valvola.

Lo slider è abilitato soltanto in modalità \code{MANUAL}, quando il CUS è disponibile.
Durante la regolazione viene mostrato il valore selezionato, mentre la richiesta di apertura viene inviata al termine della modifica.

La dashboard aggiorna i valori del sistema sulla base dello stato restituito dal backend, senza considerare automaticamente applicato un comando appena inviato.

Durante l'elaborazione delle richieste, i relativi controlli vengono temporaneamente disabilitati.

Le prove funzionali hanno verificato l'interfaccia anche durante l'indisponibilità e il ripristino del CUS.
Il DBS non verifica direttamente la disponibilità del WCS.
